\documentclass[11pt, a4paper]{article}
\usepackage[T1]{fontenc}
\usepackage[utf8]{inputenc}
\usepackage{polski}
\usepackage[margin=2cm]{geometry}
\usepackage{graphicx}
\usepackage{booktabs}
\usepackage{tabularx}
\usepackage{amsmath}
\usepackage{float}
\usepackage{titlesec}

\titlespacing{\section}{0pt}{10pt}{5pt}
\titlespacing{\subsection}{0pt}{8pt}{4pt}

\title{\vspace{-1.5cm}\textbf{Sprawozdanie \\
Weighted Maximum Coverage Problem}}
\author{Seweryn Nekrasz, Krystian Piszczek}

\begin{document}
\maketitle
\vspace{-1cm}

\section{Wstęp}
Celem projektu jest rozwiązanie wariantu problemu maksymalnego pokrycia (ang. \textit{Maximum Coverage Problem}) z uwzględnieniem wag celów oraz ograniczeń budżetowych. Zbadano i porównano trzy metody optymalizacji: solver dokładny Gurobi, Algorytm Genetyczny (AG) oraz Dyskretny Algorytm Optymalizacji Żerowania Bakterii (DBFO). 

\section{Sformułowanie problemu}
Zadanie polega na maksymalizacji sumy wag pokrytych celów, przy czym radary różnych typów mogą być rozmieszczane tylko w dozwolonych lokalizacjach, a ich sumaryczny koszt nie może przekroczyć narzuconego budżetu $B$. Przestrzeń decyzyjna ma charakter dyskretny (wybór typu radaru w danej lokalizacji).

\section{Opis implementacji}
\begin{itemize}
    \item \textbf{Gurobi:} Wykorzystany jako metoda referencyjna do wyznaczania optimum globalnego (Mixed-Integer Linear Programming).
    \item \textbf{Algorytm Genetyczny (AG):} Reprezentacja wektora dyskretnego, turniejowa selekcja, krzyżowanie jednopunktowe oraz mutacja losowa.
    \item \textbf{Dyskretne BFO (DBFO):} Modyfikacja algorytmu BFO, w której ciągła przestrzeń poszukiwań mapowana jest na k stanów dyskretnych (reprezentujących brak radaru lub wybór jednego z dostępnych typów) z wykorzystaniem funkcji sigmoidalnej i podziału przedziału prawdopodobieństwa na równe progi decyzyjne.
\end{itemize}

\section{Wyniki eksperymentów}
Badania przeprowadzono dla instancji o różnym rozmiarze (liczbie lokalizacji i celów). Tabela \ref{tab:wyniki} prezentuje uśrednione wartości zysku oraz czasu wykonania dla $N$ uruchomień heurystyk.

\begin{table}[H]
\centering
\caption{Porównanie zysku i czasu obliczeń algorytmów}
\label{tab:wyniki}
\begin{tabularx}{\textwidth}{l *{6}{X}}
\toprule
\textbf{Rozmiar} & \multicolumn{2}{c}{\textbf{Gurobi}} & \multicolumn{2}{c}{\textbf{AG}} & \multicolumn{2}{c}{\textbf{DBFO}} \\
\cmidrule(lr){2-3} \cmidrule(lr){4-5} \cmidrule(lr){6-7}
(Lok/Cele) & Zysk & Czas [s] & Śr. Zysk & Czas [s] & Śr. Zysk & Czas [s] \\
\midrule
Mały (5/7)  & 100 & 0.05 & [Brak] & [Brak] & [Brak] & [Brak] \\
Średni      & [Brak] & [Brak] & [Brak] & [Brak] & [Brak] & [Brak] \\
Duży        & [Brak] & [Brak] & [Brak] & [Brak] & [Brak] & [Brak] \\
\bottomrule
\end{tabularx}
\end{table}

\section{Analiza statystyczna}


\section{Wnioski}


\end{document}